% Generated draft: requires mathematical and source review.
\documentclass[11pt]{article}
\usepackage[margin=25mm]{geometry}
\usepackage{fontspec}
\setmainfont{Latin Modern Roman}
\usepackage{amsmath,amssymb,mathtools,unicode-math}
\setmathfont{Latin Modern Math}
\usepackage{enumitem,needspace,xurl,hyperref}
\setlength{\parindent}{0pt}
\setlength{\parskip}{5pt}
\setlength{\emergencystretch}{3em}
\newcommand{\N}{\mathbb{N}}
\newcommand{\Z}{\mathbb{Z}}
\newcommand{\Q}{\mathbb{Q}}
\newcommand{\R}{\mathbb{R}}
\newcommand{\C}{\mathbb{C}}
\newcommand{\F}{\mathbb{F}}
\newcommand{\A}{\mathbb{A}}
\newcommand{\PP}{\mathbb{P}}
\newcommand{\E}{\mathbb{E}}
\newcommand{\eps}{\varepsilon}
\newcommand{\dd}{\,\mathrm{d}}
\begin{document}
\section*{UnsolvedMath ID: 20002300}
\textbf{Source problem code:} catalogue entry
\subsection{Definitions and mathematical statement}
Miller's Question 16 asks for the Turing degrees and relative reducibilities of
\[
 \operatorname{HTP}_{\infty}(\mathbb Q),\quad
 \operatorname{HTP}_{\infty,1}(\mathbb Q),\quad
 \operatorname{HTP}_{\infty}(\mathbb Z),\quad
 \operatorname{HTP}_{\infty,1}(\mathbb Z),
\]
and records \(0'\leq_T\operatorname{HTP}_\infty(\mathbb Z)\leq_T0''\). The notation is not defined in the fragment. The aid explicitly reconstructs \(\infty,1\) as infinitely many distinct first-coordinate values, not a one-variable polynomial problem.

Fix an effective encoding of integer-coefficient polynomials with a declared finite list of variables. For \(R=\mathbb Z\) or \(\mathbb Q\), write
\[
 S_R(f)=\{\mathbf x\in R^n:f(\mathbf x)=0\}.
\]
Define \(\operatorname{HTP}(R)\) by nonemptiness of this set, \(\operatorname{HTP}_\infty(R)\) by infinitude of the whole solution set, and \(\operatorname{HTP}_{\infty,1}(R)\) by infinitude of \(\{x_1:\exists x_2,\ldots,x_n\ f(\mathbf x)=0\}\). Here \(\leq_m,\leq_{tt},\leq_T\) mean many-one, truth-table and Turing reducibility, and \(0',0''\) are the first two Turing jumps.

The aid reports
\[
 \operatorname{HTP}(R)\leq_m\operatorname{HTP}_\infty(R)
 \leq_{tt}\operatorname{HTP}_{\infty,1}(R),
\]
and both infinitude sets are \(\Pi^0_2\). The first reduction introduces a free solution coordinate; the second asks nonadaptively about each coordinate projection, since a set in a finite product is infinite exactly when some coordinate projection is infinite. Counting arbitrarily many distinct effective tuples, or first-coordinate values, gives the upper complexity bound.

Uniform MRDP is reported to yield
\[
 \operatorname{HTP}_{\infty,1}(\mathbb Z)
 \text{ is }\Pi^0_2\text{-complete},\qquad
 \operatorname{HTP}_{\infty,1}(\mathbb Z)\equiv_T0''.
\]
It represents a uniformly enumerated c.e. set as the distinguished-coordinate projection of an integer Diophantine set; infinitude of c.e. sets is \(\Pi^0_2\)-complete. Infinite auxiliary witnesses do not affect that projected infinitude, but can affect total-solution infinitude.

The aid's finite-fiber transfer criterion is: for a surjective projection \(\pi:T\to S\) preserving the distinguished coordinate, their first-coordinate projections agree regardless of fiber sizes. If in addition every fiber over a complete tuple \(s\in S\) is finite, then \(T\) is infinite exactly when \(S\) is infinite. No uniform numerical fiber bound is required, but mere surjectivity is insufficient for this latter statement.

Accordingly a uniform finite-fold MRDP hypothesis would also give \(\Pi^0_2\)-completeness of total integer-solution infinitude. Separately, an existential Diophantine definition of \(\mathbb Z\) in \(\mathbb Q\) would transfer the projected completeness result to \(\operatorname{HTP}_{\infty,1}(\mathbb Q)\), of degree \(0''\). Transferring total infinitude needs the additional finite-witness hypothesis; it is not a consequence of an unrestricted existential definition alone. These are conditional consequences, not supplied solutions of those open assumptions. The aid leaves the exact degree of \(\operatorname{HTP}_\infty(\mathbb Z)\) and both rational variants unresolved.
\subsection{Short English statement}
The supplied reconstruction separates infinitely many tuples from infinitely many first-coordinate values. Uniform MRDP settles the latter integer problem at degree \(0''\). Total infinitude is sensitive to infinitely many auxiliary witnesses, so finite-fold and rational definability consequences remain explicitly conditional; the other exact degrees are not supplied.
\subsection{Sources}
\begin{itemize}
\item Local source record: catalogue entry, catalogue entry; canonical \texttt{aim-logic-notes.json}, \texttt{notes[75]}, Question 16 attributed to Miller. Title: ``Turing degrees of infinite-solution Hilbert's Tenth Problems.''
\item AIM 2013 workshop problem list \emph{Definability and Decidability Problems in Number Theory}, supplied reference \url{https://aimath.org/pastworkshops/definabilityinntproblems.pdf}. The original displayed integer total-infinitude bounds are retained. This draft does not access a PDF.
\item Supplied research attempt 1 invokes uniform MRDP and infinitude of c.e. sets, and adds the finite-fiber transfer criterion and labeled conditional reductions. No full authoritative bibliographic entries are supplied.
\item Registry dated 14 August 2026: ``partially\_solved,'' ``PARTIAL-PROGRESS,'' statement status ``exact''; attempt partial, research-run status open, rigor ``proved,'' novelty confidence low. The machine-generated, non-peer-reviewed aid requires independent verification before citation.
\end{itemize}
\end{document}
